\section{Notation and conventions}
\label{sec:notation}

These notes are bilingual: the derivations are written in English and the
procedural summaries (\emph{resumen del proceso}, \emph{planteamiento}) are kept
in Spanish, as in the original handwritten notes.

Throughout, $\vb{X}^\tp=(X_1,X_2,\dots,X_p)$ is a random vector of $p$ observed
variables with mean vector $\mub=\E(\vb{X})$ and covariance matrix
\[
  \Si=\Cov(\vb{X})=\E\sqty{(\vb{X}-\mub)(\vb{X}-\mub)^\tp},
\]
with diagonal entries $\Si_{ii}=\si_{ii}=\Var(X_i)$.

\begin{itemize}
  \item $\la_1\ge\la_2\ge\dots\ge\la_p\ge 0$ are the eigenvalues of $\Si$
        (it is symmetric positive semidefinite, so they are real and nonnegative),
        with associated orthonormal eigenvectors $\vb{e}_1,\dots,\vb{e}_p$, i.e.
        $\vb{e}_i^\tp\vb{e}_k=\delta_{ik}$.
  \item $\rhob$ denotes the correlation matrix and $\vb{Z}$ the vector of
        standardized variables.
  \item Transposition is written $\vb{a}^\tp$; in the original notes this appears
        both as $\vb{a}'$ and as $\mat{Q}^{T}$, and both mean the same thing here.
  \item $m$ always denotes the number of common factors in
        Section~\ref{sec:factor_analysis}, with $m<p$.
  \item $G$ denotes the categorical class label and $K$ the number of classes in
        Section~\ref{sec:discriminant_analysis}; the handwritten notes use both
        $N$ and $K$ for this count, and here it is always $K$.
  \item In Section~\ref{sec:logistic_regression}, $\si(\cdot)$ denotes the
        logistic function $\si(z)=1/(1+e^{-z})$ --- not to be confused with the
        entries $\si_{ii}$ of $\Si$ --- and $\ell(\be)$ the log-likelihood.
        The handwritten notes write the logistic function as $\ln(z)$, which
        here is reserved for the natural logarithm.
\end{itemize}
